\newpage
\section{第24章：存在类型}

\begin{taplauthorsolutionentrybox}
\phantomsection\label{sol:author:24.1.1}
\textbf{24.1.1 解答}

\(p_6\) 同时提供常量 \(a\) 和函数 \(f\)，但客户只能反复把 \(f\) 应用于
\(a\)，最后仍得到无法观察的抽象 \(X\) 值。\(p_7\) 允许用 \(f\) 构造
\(X\) 值，却仍没有观察它们的操作。\(p_8\) 的两个字段都可使用，但接口没有
隐藏任何与表示有关的信息，存在包也就失去作用。
\taplsolutionbacklink{ex:24.1.1}{返回练习24.1.1}
\end{taplauthorsolutionentrybox}

\begin{taplauthorsolutionentrybox}
\phantomsection\label{sol:author:24.2.1}
\textbf{24.2.1 解答}
\[
\begin{aligned}
\taplterm{stackADT}={}&
\{*\tapltype{List}\,\tapltype{Nat},
\{\taplterm{new}=\taplterm{nil}[\tapltype{Nat}],\\
&\taplterm{push}=\lambda n:\tapltype{Nat}.
 \lambda s:\tapltype{List}\,\tapltype{Nat}.
 \taplterm{cons}[\tapltype{Nat}]n\,s,\\
&\taplterm{top}=\lambda s:\tapltype{List}\,\tapltype{Nat}.
 \taplterm{head}[\tapltype{Nat}]s,\\
&\taplterm{pop}=\lambda s:\tapltype{List}\,\tapltype{Nat}.
 \taplterm{tail}[\tapltype{Nat}]s,\\
&\taplterm{isempty}=\lambda s:\tapltype{List}\,\tapltype{Nat}.
 \taplterm{isnil}[\tapltype{Nat}]s\}\}\\
&\textbf{as }\{\exists\ \tapltype{Stack},\\
&\quad\{\taplterm{new}:\tapltype{Stack},\\
&\hspace{15mm}\taplterm{push}:\tapltype{Nat}\to\tapltype{Stack}\to\tapltype{Stack},\\
&\hspace{15mm}\taplterm{top}:\tapltype{Stack}\to\tapltype{Nat},\\
&\hspace{15mm}\taplterm{pop}:\tapltype{Stack}\to\tapltype{Stack},\\
&\hspace{15mm}
\taplterm{isempty}:\tapltype{Stack}\to\tapltype{Bool}\}\}.
\end{aligned}
\]
检查器给出
\[
\begin{aligned}
?\ \taplterm{stackADT}:\{\exists\ \tapltype{Stack},
\{&\taplterm{new}:\tapltype{Stack},
\taplterm{push}:\tapltype{Nat}\to\tapltype{Stack}\to\tapltype{Stack},\\
&\taplterm{top}:\tapltype{Stack}\to\tapltype{Nat},
\taplterm{pop}:\tapltype{Stack}\to\tapltype{Stack},\\
&\taplterm{isempty}:\tapltype{Stack}\to\tapltype{Bool}\}\}
\end{aligned}
\]
打开包后，
\(\taplterm{stack.top}(\taplterm{stack.push}\ 5
(\taplterm{stack.push}\ 3\ \taplterm{stack.new}))\) 得到 \(5\)。
\taplsolutionbacklink{ex:24.2.1}{返回练习24.2.1}
\end{taplauthorsolutionentrybox}

\begin{taplauthorsolutionentrybox}
\phantomsection\label{sol:author:24.2.2}
\textbf{24.2.2 解答}
\[
\begin{aligned}
\taplterm{counterADT}={}&
\{*\tapltype{Ref}\,\tapltype{Nat},
\{\taplterm{new}=\lambda\_:\tapltype{Unit}.\taplterm{ref}\ 1,\\
&\taplterm{get}=\lambda r:\tapltype{Ref}\,\tapltype{Nat}.\,!r,\\
&\taplterm{inc}=\lambda r:\tapltype{Ref}\,\tapltype{Nat}.
  r:=\taplterm{succ}(!r)\}\}\\
&\textbf{as }\{\exists\ \tapltype{Counter},\\
&\quad\{\taplterm{new}:\tapltype{Unit}\to\tapltype{Counter},
\taplterm{get}:\tapltype{Counter}\to\tapltype{Nat},
\taplterm{inc}:\tapltype{Counter}\to\tapltype{Unit}\}\}.
\end{aligned}
\]
检查器给出
\[
\begin{aligned}
?\ \taplterm{counterADT}:{}&
\{\exists\ \tapltype{Counter},\\
&\{\taplterm{new}:\tapltype{Unit}\to\tapltype{Counter},\\
&\quad\taplterm{get}:\tapltype{Counter}\to\tapltype{Nat},\\
&\quad\taplterm{inc}:\tapltype{Counter}\to\tapltype{Unit}\}\}
\end{aligned}
\]
\taplsolutionbacklink{ex:24.2.2}{返回练习24.2.2}
\end{taplauthorsolutionentrybox}

\begin{taplauthorsolutionentrybox}
\phantomsection\label{sol:author:24.2.3}
\textbf{24.2.3 解答}

令 \(\taplterm{zeroCounter}\) 为零值计数器对象，则可取
\[
\begin{aligned}
\tapltype{FlipFlop}
={}&\{\exists X,\{\taplterm{state}:X,\\
&\quad\taplterm{methods}:\{
\taplterm{read}:X\to\tapltype{Bool},\\
&\hspace{34mm}
\taplterm{toggle}:X\to X,
\taplterm{reset}:X\to X\}\}\},\\
\taplterm{f}
={}&\{*\tapltype{Counter},
\{\taplterm{state}=\taplterm{zeroCounter},\\
&\quad\taplterm{methods}=\{
\taplterm{read}=\lambda s:\tapltype{Counter}.
\taplterm{iseven}(\taplterm{sendget}\ s),\\
&\quad\taplterm{toggle}=\lambda s:\tapltype{Counter}.
\taplterm{sendinc}\ s,\,
\\
&\quad
\taplterm{reset}=\lambda s:\tapltype{Counter}.
\taplterm{zeroCounter}\}\}\}\\
&\textbf{as }\tapltype{FlipFlop}.
\end{aligned}
\]
检查器给出 \(?\ \taplterm{f}:\tapltype{FlipFlop}\)。
\taplsolutionbacklink{ex:24.2.3}{返回练习24.2.3}
\end{taplauthorsolutionentrybox}

\begin{taplauthorsolutionentrybox}
\phantomsection\label{sol:author:24.2.4}
\textbf{24.2.4 解答}
\[
\begin{aligned}
c=\{*\tapltype{Ref}\,\tapltype{Nat},
\{&\taplterm{state}=\taplterm{ref}\ 5,\\
&\taplterm{methods}=\{
\taplterm{get}=\lambda x:\tapltype{Ref}\,\tapltype{Nat}.\,!x,\\
&\taplterm{inc}=\lambda x:\tapltype{Ref}\,\tapltype{Nat}.
(x:=\taplterm{succ}(!x);x)\}\}\}
\textbf{ as }\tapltype{Counter}.
\end{aligned}
\]
\taplsolutionbacklink{ex:24.2.4}{返回练习24.2.4}
\end{taplauthorsolutionentrybox}

\begin{taplauthorsolutionentrybox}
\phantomsection\label{sol:author:24.2.5}
\textbf{24.2.5 解答}

这一类型确实能在单个对象内部实现并集，却几乎无法调用。要调用某对象的
\(\taplterm{union}:X\to X\to X\)，必须提供两个同一表示类型 \(X\) 的状态。
打开两个不同对象会分别引入两个互不相同的抽象类型，第二个对象的状态不能传给
第一个对象的方法；允许这样传递也不安全，因为二者表示可能完全不同。因此，
该接口至多允许一个集合与自身求并。
\taplsolutionbacklink{ex:24.2.5}{返回练习24.2.5}
\end{taplauthorsolutionentrybox}

\begin{taplauthorsolutionentrybox}
\phantomsection\label{sol:author:24.3.2}
\textbf{24.3.2 解答}

至少要证明翻译保持赋型和计算。若以 \(\llbracket-\rrbracket\) 表示翻译，
则应有
\[
\Gamma\vdash t:T
\Longrightarrow
\llbracket\Gamma\rrbracket\vdash
\llbracket t\rrbracket:\llbracket T\rrbracket
\]
以及
\[
t\longrightarrow^*t'
\Longrightarrow
\llbracket t\rrbracket\longrightarrow^*\llbracket t'\rrbracket.
\]
这两项容易验证。反向性质则不成立：翻译可能把源语言中错误或卡住的项映射为
目标语言中良类型且可继续求值的项。具体反例是
\[
(\{*\tapltype{Nat},0\}\textbf{ as }\{\exists X,X\})[\tapltype{Bool}].
\]
源项试图把存在包当作全称值作类型应用，因而既非良类型也无法求值；编码后的
目标项却是良类型的，而且不会卡住。
\taplsolutionbacklink{ex:24.3.2}{返回练习24.3.2}
\end{taplauthorsolutionentrybox}

\begin{taplauthorsolutionentrybox}
\phantomsection\label{sol:author:24.3.3}
\textbf{24.3.3 解答}

作者没有见过这一编码的完整书面版本。它似乎可以做到，但不会是逐个局部构造
展开的简单句法糖；翻译必须一次处理整个程序。
\taplsolutionbacklink{ex:24.3.3}{返回练习24.3.3}
\end{taplauthorsolutionentrybox}
